\chapter{CI/CD --- GitHub Actions}

\section{Presentation}

Le code source est heberge sur le depot GitHub
\texttt{Michael-Rousseau/devSecOps}\footnote{\texttt{https://github.com/Michael-Rousseau/devSecOps}}.
Quatre workflows GitHub Actions automatisent l'ensemble du cycle de vie~:

\begin{table}[H]
    \centering
    \begin{tabularx}{\textwidth}{lllX}
        \toprule
        \textbf{Workflow} & \textbf{Declencheur} & \textbf{Type} & \textbf{Actions} \\
        \midrule
        ci.yml & PR vers main & CI & Lint, test, build (frontend + backend) \\
        cd-frontend.yml & Push sur main & CD & Build Angular, sync S3 (\texttt{deus-dashboard-frontend-dev}) \\
        cd-backend.yml & Push sur main & CD & Build Docker, push ECR, deploy ECS \\
        terraform.yml & PR / Push & IaC & Plan sur PR, Apply sur merge \\
        \bottomrule
    \end{tabularx}
    \caption{Workflows GitHub Actions}
\end{table}

\section{Pipeline CI}

Le workflow CI s'execute sur chaque Pull Request~:

\begin{lstlisting}[language=bash,caption={Jobs CI en parallele}]
frontend:
  npm ci -> npm run lint -> npm run build

backend:
  cd backend && npm ci -> npm run build
\end{lstlisting}

La concurrence est geree pour annuler les executions obsoletes sur une meme branche.

\section{Deploiement Frontend}

\begin{enumerate}
    \item Build de l'application Angular en mode production
    \item Synchronisation des fichiers vers le bucket S3
        \texttt{deus-dashboard-frontend-dev} via \texttt{aws s3 sync}
\end{enumerate}

\textbf{Note~:} L'etape d'invalidation CloudFront initialement prevue a ete
supprimee car l'hebergement utilise le mode S3 Static Website Hosting
(cf. chapitre~\ref{chap:architecture} --- Adaptations Learner Lab).

\section{Deploiement Backend}

\begin{enumerate}
    \item Build de l'image Docker
    \item Push vers Amazon ECR
        \texttt{459624991030.dkr.ecr.us-east-1.amazonaws.com/deus-dashboard-backend}
        (tag~: SHA du commit + latest)
    \item Mise a jour de la Task Definition ECS sur le cluster
        \texttt{deus-dashboard-dev}
    \item Deploiement du service ECS \texttt{deus-dashboard-backend}
        avec attente de stabilite
\end{enumerate}

\section{Terraform en CI}

\begin{description}
    \item[Sur Pull Request] \texttt{terraform plan} est execute et le resultat
        est poste en commentaire sur la PR pour revue.
    \item[Sur merge vers main] \texttt{terraform apply -auto-approve} est execute
        automatiquement.
\end{description}

\section{Gestion des credentials AWS Learner Lab}

Les credentials AWS Learner Lab (compte \texttt{459624991030}, region
\texttt{us-east-1}) expirent toutes les \~{}4 heures.
Un script \texttt{scripts/update-aws-keys.sh} permet de mettre a jour
simultanement~:
\begin{itemize}
    \item \texttt{\~{}/.aws/credentials} (CLI locale)
    \item Les GitHub Secrets du depot \texttt{Michael-Rousseau/devSecOps}
        (via \texttt{gh secret set})
\end{itemize}

Les trois secrets concernes sont~:
\begin{itemize}
    \item \texttt{AWS\_ACCESS\_KEY\_ID}
    \item \texttt{AWS\_SECRET\_ACCESS\_KEY}
    \item \texttt{AWS\_SESSION\_TOKEN}
\end{itemize}

\textbf{Limitation~:} L'absence de federation OIDC (non autorisee par Learner Lab)
impose l'utilisation de cles statiques dans les GitHub Secrets. Ces cles doivent
etre renouvelees manuellement avant chaque session de deploiement. En production,
on privilegierait un provider OIDC
(\texttt{aws-actions/configure-aws-credentials}) pour eviter le stockage de cles
longue duree.

\begin{figure}[H]
    \centering
    \includegraphics[width=\textwidth]{gh_watch_ci.png}
    \caption{Pipeline CI sur une Pull Request}
\end{figure}
